\section{Model and estimation specification}
\label{app:spec}

This appendix is generated from the code that ran: the constants, bounds, solver settings and loss terms are
read out of the modules that execute during inference (\texttt{evaluation/export\_spec.py}), so it cannot
disagree with what was computed.

\subsection{State, parameters and constants}

The engine's state vector has twenty-one entries, of which eight are evolved in this stage: glucose,
insulin, insulin action, stomach and gut glucose, glycogen, ketones and cumulative energy expenditure. The
postprandial dynamics of glucose are carried by the first five (Section~\ref{sec:model}); the others
enter the glucose equation only through terms that are numerically negligible for the initial state used
(a glycogen level far above the threshold of the ketone drive).

\begin{table}[h]
\centering\small
% Generated by evaluation/export_spec.py -- DO NOT EDIT.
\begin{tabular}{lrrl}
\toprule
Parameter & Population default & Inference bounds & Status \\
\midrule
$S_I$ & 1 & $[0.3, 1.6]$ & free \\
$k_e$ & 0.026 & $[0.015, 0.04]$ & free \\
$k_a$ & 0.022 & $[0.012, 0.032]$ & free \\
$W$ & 78 & -- & fixed \\
$c_{\mathrm{circ}}$ & 0 & -- & fixed \\
$\mathrm{RMR}$ & 1.06 & -- & fixed \\
$p_1$ & 0.026 & -- & fixed \\
$p_2$ & 0.028 & -- & fixed \\
$p_3$ & 1.3e-05 & -- & fixed \\
$n$ & 0.16 & -- & fixed \\
$\beta$ & 0.16 & -- & fixed \\
\bottomrule
\end{tabular}

\begin{tabular}{lr}
\toprule
Constant & Value \\
\midrule
\texttt{basal\_glucose\_mg\_dl} & 90 \\
\texttt{basal\_insulin\_uU\_ml} & 10 \\
\texttt{glucose\_distribution\_volume\_per\_kg\_dl} & 1.6 \\
\texttt{hepatic\_cortisol\_threshold\_ug\_dl} & 15 \\
\texttt{hepatic\_gain\_mg\_dl\_min\_per\_ug\_dl} & 0.1 \\
\texttt{glycogen\_restoration\_insulin\_threshold\_uU\_ml} & 15 \\
\texttt{glycogen\_restoration\_rate\_g\_min} & 0.15 \\
\texttt{ketone\_drive\_rate\_mmol\_l\_min} & 0.12 \\
\texttt{ketone\_decay\_per\_min} & 0.02 \\
\texttt{gluconeogenesis\_threshold\_mg\_dl} & 72 \\
\texttt{carb\_to\_glucose\_fraction} & 0.9 \\
\texttt{meal\_pulse\_sigma\_min} & 2 \\
\bottomrule
\end{tabular}

\caption{Left: the free parameters (rows one to three) and the fixed parameters of the model, with
population values and inference bounds. Right: the fixed constants.}
\label{tab:spec}
\end{table}

\paragraph{Meal input.}
A meal of $c$ grams of carbohydrate, $f$ of fat and $b$ of fibre delivers a glucose mass
$M=0.9\cdot 1000\,c\,/(1+0.08\,b+0.005\,f)$\,mg as a Gaussian rate pulse of width $2$\,min centred on the
meal, integrating to $M$. The simulation lasts $210$\,min with the meal at $30$\,min and is saved every
$5$\,min.

\paragraph{Solver.}
Tsit5 with a PID step-size controller (relative and absolute tolerance $10^{-4}$, at most $10^{4}$ steps),
differentiated by reverse mode through the solver, in double precision.

\subsection{Observables and loss}

With the baseline $\bar G_{\mathrm{pre}}$ the mean of the six samples before the meal, and the smooth
positive part $[x]_+^{\beta}=\beta^{-1}\log(1+e^{\beta x})$ with $\beta=10$,
\begin{align}
\iauc&=\int_{0}^{T}[G-\bar G_{\mathrm{pre}}]_+^{\beta}\,dt \quad (\text{trapezoid on the sample grid}),\\
\text{centroid}&=\int_0^T t\,[G-\bar G_{\mathrm{pre}}]_+^{\beta}\,dt\Big/\iauc,
\end{align}
and the trace is the vector of $[G-\bar G_{\mathrm{pre}}]_+^{\beta}$ at the samples of the window, of
length $T=180$\,min for CGMacros and Shanghai and $145$\,min for the cohort with standardized meals. The
Shanghai trace uses every third sample, the real $15$-minute measurements.

For the prediction fits the loss is the masked mean squared error of the iAUC over the training meals plus
$\lambda\sum_j\bigl((\theta_j-\theta_j^{\mathrm{pop}})/\mathrm{range}_j\bigr)^2$ with $\lambda=0.01$,
minimized by Adam (learning rate $0.02$, $150$ steps) with projection onto the box after each update. For
every identifiability analysis $\lambda=0$ and the loss is the Gaussian negative log-likelihood
$\tfrac12\sum_k r_k^2/\sigma_k^2$ over residuals $r_k$ of the stacked, weighted and (for the trace) whitened
observables, with frozen $\sigma_k$ and, for the trace, frozen autocorrelation $\rho$.

\subsection{Profile likelihood}

For each parameter the grid has $15$ log-spaced points over the box. At every point the other parameters
are re-optimized with $100$ Adam steps in log coordinates started from the estimate (a warm start that keeps
the inner optimum on one branch). The reference level is the lower of the loss at the estimate and the lowest
profile value, so an estimate that is not an optimum is flagged and not silently moved. A grid point whose
inner loss fell by more than $0.05$ over the last quarter of the steps is counted as unsettled and reported
(it overstates the profile there, which flatters identifiability). For $\ke$ and $\ka$ six further points
extend the grid to four times the upper bound; they enter no classification.

\subsection{Coordinates}

In $(\log\SI,\log\tau_1,\log p)$ the box is the rectangular hull of the image of the rate box (a superset,
deliberately: a tighter box would cut into the likelihood and manufacture identifiability), and the real-pole
condition $\log p\le2\log\tau_1-\log 4$ is enforced by projection onto the tied-rate boundary $\ke=\ka$. Every
coordinate fit records whether its implied rates fall inside the original box and whether it is on the
boundary. The canonical engine integrates directly in $(\ke+\ka,\ke\ka)$ and never takes a square root of
the rates.

\subsection{Replica and synthetic study}

The replica noise is first-order autoregressive with the lag-one correlation and innovation standard
deviation of the residual of the real maximum-likelihood fit; when that estimate is unstable (innovation
standard deviation outside $0.5$ to $60$\,mg\,dL$^{-1}$) the fallback is a marginal standard deviation of
$10$\,mg\,dL$^{-1}$ and a correlation of $0.7$. Carbohydrate error is $\log$-normal,
$\exp(\mathcal N(-\sigma^2/2,\sigma^2))$ with $\sigma^2=\log(1+\mathrm{CV}^2)$, so it has unit mean. The
synthetic study draws $\SI$, $\ke$ and $\ka$ uniformly and independently over the box, at least $5\%$ of its
width from either bound, with a stream seeded by the subject and the seed.

\subsection{Second model class}

The Dalla Man model (three-compartment gut with a $\tanh$ gastric-emptying law, two glucose and two insulin
compartments, delayed insulin action on endogenous production, a subcutaneous glucose compartment) is run in
double precision with relative and absolute tolerance $10^{-6}$. The six fitted parameters and their boxes
are in Table~\ref{tab:dm}; the basal glucose is the subject's mean pre-meal CGM value and the insulin
secretion slope is fixed.

\begin{table}[h]
\centering\small
\begin{tabular}{llc}
\toprule
parameter & meaning & box \\
\midrule
$V_{mx}$ & insulin sensitivity & $[0.005,\,0.150]$ \\
$k_{abs}$ & intestinal absorption & $[0.020,\,0.200]$ \\
$k_{max}$ & maximal gastric emptying & $[0.020,\,0.120]$ \\
$k_{min}$ & minimal gastric emptying & $[0.005,\,0.080]$ \\
$f$ & fraction absorbed & $[0.60,\,1.00]$ \\
$T_d$ & sensor lag (min) & $[5,\,25]$ \\
\bottomrule
\end{tabular}
\caption{Parameters fitted in the second model class and their boxes (log scale for fitting).}
\label{tab:dm}
\end{table}
